\chapter{Methodology}\label{chap:methods}

This chapter describes the detector (Sections~\ref{sec:overview}--\ref{sec:fusion}), the synthetic training data tested for its judge (Section~\ref{sec:synth}), and the protocol under which every design choice was made (Section~\ref{sec:protocol}). Figure~\ref{fig:arch} shows the whole system.

\section{Overview of the detector}\label{sec:overview}

The detector has two branches that read the same record $x = (I, q, r)$, and one fusion step.

\begin{itemize}
    \item \textbf{Generative branch.} A VLM judge, fine-tuned to list the hallucinated spans of $r$ as JSON, is decoded once greedily and $N = 8$ times by sampling. A label-aware span decoder converts each decode into character-offset spans. Together they give nine sets of labelled spans.
    \item \textbf{Uncertainty branch.} A second copy of the same backbone, frozen and \emph{not} fine-tuned, reads $r$ as if it were its own answer to $q$ about $I$. From this pass we take four signals at every answer token: the entropy of the next-token distribution, the negative log-likelihood of the observed token, the share of attention each head pays to the image (look-back ratio), and the last-layer hidden state. A BiGRU encoder with five binary CRF heads decodes them into a tenth set of labelled spans.
    \item \textbf{Fusion.} A per-label, per-character vote keeps a character for category $L$ when at least $k = 3$ of the ten voters mark it with $L$. Each kept character gets its own probability, the mean of the voters' probabilities for that label. Consecutive characters with the same label and the same rounded probability form one output span.
\end{itemize}

The two branches are deliberately different in kind. The generative branch is supervised and black-box: it only needs the judge's output text. The uncertainty branch reads model internals and is trained separately on the same data. Each component is motivated by an experiment in Chapter~\ref{chap:experiments}: the backbone by the Part~1 comparison, the voter and the probe by the Part~3 ablations, and the graded probability by the analysis of the Cor metric.

\begin{figure}[!htb]
\centering
\resizebox{\textwidth}{!}{%
\begin{tikzpicture}[
    font=\small,
    box/.style={draw, rounded corners=2pt, align=center, minimum height=1.25cm, minimum width=3.2cm, inner sep=4pt, fill=white},
    data/.style={draw, align=center, minimum height=1.0cm, inner sep=3pt, fill=gray!8},
    gen/.style={box, fill=blue!7},
    unc/.style={box, fill=orange!10},
    fus/.style={box, fill=green!9},
    arr/.style={-{Stealth[length=2.2mm]}, thick}]
  % input
  \node[data, minimum width=2.4cm] (in) at (0, 0) {Record $x$:\\image $I$,\\question $q$,\\answer $r$};
  % generative branch (top row)
  \node[gen] (judge)  at (3.7,  1.6) {\textbf{A. Fine-tuned judge}\\Qwen3-VL-8B\\+ QLoRA};
  \node[gen] (dec)    at (7.4,  1.6) {1 greedy decode\\+ 8 samples\\($T = 0.7$)};
  \node[gen] (spanio) at (11.1, 1.6) {\textbf{B. Span decoder}\\JSON text $\rightarrow$\\char offsets};
  % uncertainty branch (bottom row)
  \node[unc] (frozen) at (3.7,  -1.6) {\textbf{D. Frozen judge}\\Qwen3-VL-8B,\\teacher-forced on $r$};
  \node[unc] (feat)   at (7.4,  -1.6) {token features:\\entropy, NLL,\\look-back, hidden};
  \node[unc] (crf)    at (11.1, -1.6) {BiGRU encoder\\+ 5 binary CRFs\\(probe)};
  % fusion (right column)
  \node[fus] (vote)  at (15.6,  1.6) {\textbf{C. Per-label vote}\\10 voters,\\keep $\mathrm{votes}_L(c) \ge 3$};
  \node[fus] (grade) at (15.6, -1.6) {\textbf{E. Graded output}\\per-character\\probability $\hat p_L(c)$};
  \node[data, minimum width=2.9cm] (out) at (15.6, -3.75) {labelled spans $(a, b, p, L)$};
  % arrows
  \draw[arr] (in.east) -- ++(0.35,0) |- (judge.west);
  \draw[arr] (in.east) -- ++(0.35,0) |- (frozen.west);
  \draw[arr] (judge) -- (dec);
  \draw[arr] (dec) -- (spanio);
  \draw[arr] (frozen) -- (feat);
  \draw[arr] (feat) -- (crf);
  \draw[arr] (spanio) -- node[above, font=\footnotesize] {9 voters} (vote);
  \draw[arr] (crf.east) -- ++(0.5,0) |- node[pos=0.25, right, font=\footnotesize] {1 voter} ([yshift=-3mm]vote.west);
  \draw[arr] (vote) -- (grade);
  \draw[arr] (grade) -- (out);
  % branch frames
  \begin{scope}[on background layer]
    \node[draw=blue!45, dashed, rounded corners, fit=(judge)(dec)(spanio), inner sep=6pt,
          label={[font=\footnotesize\itshape, text=blue!60!black]above:generative branch (supervised, black-box)}] {};
    \node[draw=orange!60, dashed, rounded corners, fit=(frozen)(feat)(crf), inner sep=6pt,
          label={[font=\footnotesize\itshape, text=orange!60!black]below:uncertainty branch (white-box probe)}] {};
  \end{scope}
\end{tikzpicture}%
}
\caption{Architecture of the detector. The generative branch (A, B) and the uncertainty branch (D) each produce labelled spans; the fusion step (C, E) votes on characters per label and keeps a graded per-character probability.}
\label{fig:arch}
\end{figure}

\section{Component A: the generative judge}\label{sec:judge}

\paragraph{Input and output format.}
The judge is a VLM fine-tuned as a conditional generator. Its user turn contains the image, an instruction that defines the five categories, the question $q$, and the answer $r$ to analyse. The full prompt is in \hyperref[app:prompt]{Appendix~A}. Its assistant turn is a JSON array that lists the hallucinated spans of $r$ from left to right:
\begin{center}
\texttt{[\{"text": "<verbatim span>", "prob": 0.67, "label": "invention"\}, ...]}
\end{center}
A clean answer has target \texttt{[]}. Training targets are built from the gold spans. Each span is copied verbatim from $r$ with its gold probability, so the judge learns to write down both the location (as text) and the annotator agreement (as a number). Writing spans as text rather than offsets avoids asking a language model to count characters, and it works the same way for all four languages and scripts. The cost is that its output must be mapped back to offsets (Component B).

\paragraph{Training.}
The judge is trained with QLoRA~\cite{dettmers2023qlora}:
\begin{itemize}
    \item The backbone is loaded in 4-bit NF4 and frozen.
    \item Rank-16 LoRA adapters ($\alpha = 16$)~\cite{hu2022lora} are added to the attention and MLP projections of the language model. The vision encoder is not adapted.
    \item The loss is the token-level cross-entropy on the assistant turn only, so the prompt, image and answer tokens are masked.
    \item AdamW~\cite{loshchilov2019adamw}, learning rate $2\cdot 10^{-4}$, 50 warm-up steps, then linear decay; batch size 1 with 8 gradient-accumulation steps, 2 epochs, sequences of up to 4,096 tokens.
\end{itemize}
This is the recipe of our shared-task system. It was fixed before any experiment in this thesis and never tuned on validation. On the 12,082 training records it gives 3,020 optimizer steps and takes about 15 hours on one RTX~3090 (2.2~s per example). The backbone is chosen among four VLM families (Section~\ref{sec:exp-part1}); Qwen3-VL-8B~\cite{bai2025qwen3} wins.

\section{Component B: label-aware span decoding}\label{sec:decoder}

The judge writes spans as text, and they must be placed in $r$. A span text $t$ can occur several times in $r$, and the judge reproduces it with small deviations, for example in whitespace, punctuation or casing. The decoder therefore searches left to right with a cursor. The spans are listed in reading order, so the search for span $j$ starts after span $j-1$, with a tolerant fallback over the whole answer.

\paragraph{The dual-label problem.}
The decoder of our shared-task system, and of the common code base, dropped any span that overlapped an already placed span by more than half, \emph{whatever its label}. Gold annotations, however, often give the same text two categories (Chapter~\ref{chap:problem}). When the judge correctly wrote ``\texttt{red car}'' once as invention and once as mischaracterization, the second span was discarded. The loss is large even for perfect output. Fed the gold targets of the training records, the old decoder failed to reproduce 3,939 of the 8,646 labelled records exactly, and on the validation selection half it capped Cor+Lbl at 0.925.

\paragraph{The fix.}
In the new decoder a span clashes only with spans of its \emph{own} label; gold never overlaps two spans of the same label. A span whose label differs from that of the previous span is searched from the previous span's \emph{start}, so partly overlapping dual-label spans are found. The variant was chosen on the training round trip only, never on validation. It reduces the unreproduced records from 3,939 to 2,531, most of the rest being whitespace trimming, and lifts the perfect-output ceiling on validation from 0.925 to 0.984. Raw generations are saved, so any decoder can be re-applied without the GPU.

\section{Component C: per-label self-consistency vote}\label{sec:vote}

\paragraph{Voters.}
For each record, the judge is decoded once greedily and $N = 8$ times with temperature $T = 0.7$, and each decode goes through Component B. A single sampled decode is worse than the greedy one (mean validation Cor+Lbl 0.408 against 0.427), but where the decodes \emph{agree} the evidence is stronger. Before the vote was used, a pre-registered diversity check confirmed that the samples are not near-copies of each other: on 50 records their mean pairwise character-set Jaccard similarity was 0.71--0.77, below the threshold of 0.95.

\paragraph{Vote.}
Let $S_1, \dots, S_M$ be the span sets of the $M$ voters, with weights $w_i = 1$, and let $\mathbb{1}_{i,L}(c)$ indicate that voter $i$ covers character $c$ with a span of category $L$, and $p_{i,L}(c)$ the probability of that span (the maximum if two of its spans overlap). For every category $L$ and character $c$:
\begin{align}
    \mathrm{votes}_L(c) &= \sum_{i=1}^{M} w_i\, \mathbb{1}_{i,L}(c), \label{eq:votes}\\
    \bar p_L(c) &= \frac{1}{\sum_i w_i} \sum_{i=1}^{M} w_i\, \mathbb{1}_{i,L}(c)\, p_{i,L}(c). \label{eq:pbar}
\end{align}
Character $c$ is kept for category $L$ if $\mathrm{votes}_L(c) \ge k$. The vote is taken \emph{per label}. A character may therefore be kept for two categories, as in the gold data, and a category survives only where at least $k$ voters give that category. An earlier single-label version kept, at each character, only the label with the largest summed confidence, and so could never produce dual-label output. The threshold $k$ was chosen once among $\{2, 3, 4\}$ on the selection half and reused for all seeds: $k = 3$. In the plain voter (system F2 in Chapter~\ref{chap:experiments}) the voters are the 8 samples and the greedy decode ($M = 9$), and a run of kept characters becomes one span with probability equal to the mean of $\bar p_L$ over the run.

\section{Component D: the uncertainty probe}\label{sec:probe}

\paragraph{Frozen judge and features.}
The probe reads a frozen, \emph{not} fine-tuned Qwen3-VL-8B. It is given $(I, q)$ as a normal user turn and $r$ as its own assistant answer under teacher forcing, so the features measure how surprising $r$ is to a plain VQA model, not how well a model can list spans. For each answer token $t$ we record:
\begin{itemize}
    \item the entropy $H_t = -\sum_v \pi_t(v) \log \pi_t(v)$ of the next-token distribution $\pi_t$;
    \item the negative log-likelihood $-\log \pi_t(r_t)$ of the observed token;
    \item the look-back ratio of every attention head~\cite{chuang2024lookback}, adapted to images~\cite{huang2024opera}: the share of the head's attention mass on image tokens, relative to image plus earlier answer tokens. With 36 layers and 32 heads this gives 1,152 values;
    \item the last-layer hidden state (4,096 values).
\end{itemize}
Attention maps are reduced to ratios one layer at a time inside forward hooks, so full attention tensors are never held in memory. The two wide blocks are reduced by PCA, fitted on a 200k-token sample of training data only: the hidden state to 256 components and the look-back ratios to 128. All features are standardised on train, giving 386 features per token.

\paragraph{Encoder and CRF heads.}
The token features go through a linear layer of width 512, a bidirectional GRU~\cite{cho2014gru} with 256 units, to whose input an 8-dimensional language embedding is added, and dropout 0.2. For each category $L$ a linear layer gives emission scores $e_{t,L} \in \mathbb{R}^2$ for the two states \textsf{OUT}/\textsf{IN}, and a 2-state linear-chain CRF~\cite{lafferty2001crf} with transition matrix $A_L \in \mathbb{R}^{2 \times 2}$ and start and end scores, scores a label path $z_{1:T}$:
\begin{equation}
    \mathrm{score}_L(z) = \pi^{\mathrm{start}}_{L, z_1} + \sum_{t=1}^{T} e_{t,L}[z_t] + \sum_{t=2}^{T} A_L[z_{t-1}, z_t] + \pi^{\mathrm{end}}_{L, z_T},
    \qquad
    P_L(z \mid x) = \frac{\exp \mathrm{score}_L(z)}{\sum_{z'} \exp \mathrm{score}_L(z')}.
\end{equation}
The five CRFs share the encoder and are trained jointly on the exact sequence log-likelihood, $\mathcal{J} = -\sum_L \log P_L(z^\star_L \mid x)$, with the partition function computed by the forward algorithm. The CRF learns that spans are contiguous, so no per-token threshold has to be tuned. All five CRFs run as one batched module, checked against brute-force enumeration on toy sequences.

\paragraph{Multi-label targets.}
Token $t$ is \textsf{IN} for category $L$ if \emph{any} gold span of category $L$ covers it. Text that gold labels as both invention and mischaracterization therefore trains both heads. An earlier version gave each token only the label of the first gold span covering it, which taught the probe never to output dual labels. At inference each head is decoded independently by Viterbi. Consecutive \textsf{IN} tokens become a span, character offsets come from the tokenizer, and the span probability is the mean posterior marginal $P_L(z_t = \textsf{IN} \mid x)$ over its tokens.

\paragraph{Training.}
Adam, learning rate $10^{-3}$, batch size 16, at most 20 epochs. Training stops early on the selection-half \emph{loss}, not on the metric, with patience 3. Records are cut at 1,024 tokens. The features cover every train, validation and test record, and the evaluation code checks this: a record without features would otherwise receive an empty prediction, and for a clean record that scores a perfect 1.0 (Section~\ref{sec:exp-fixes}). The probe is trained with three seeds.

\section{Fusion and graded probabilities}\label{sec:fusion}

\paragraph{Probe as a voter.}
The probe enters the vote of Section~\ref{sec:vote} as \emph{one} more voter with the same weight as a single decode, giving $M = 10$ voters. The threshold stays at $k = 3$, the value already chosen, and is not re-tuned. This is system F3. Because the probe counts like one decode, it cannot create spans on its own. It can only add the deciding vote where two decodes already agree, or reinforce where more do.

\paragraph{Graded per-character probability.}
The span-mean probability of the plain voter gives every character of a span the same value, which the rank-based Cor and Cor+Lbl treat as ties. Gold probabilities, however, vary inside spans: the core of a hallucination is marked by more annotators than its edges. The final detector (G1) therefore keeps the per-character value
\begin{equation}
    \hat p_L(c) = \mathrm{round}\big(\bar p_L(c), 2\big) \qquad \text{for every kept } c,
    \label{eq:graded}
\end{equation}
and outputs maximal runs of characters with the same category and the same $\hat p_L$ as spans. The characters kept, and therefore IoU, are exactly those of F3. Only the probabilities change. Since $\bar p_L(c)$ grows with the number of voters that mark $c$ and with their confidence, characters on which all decodes agree rank above characters that only a few voters mark. An alternative, the plain vote fraction $\mathrm{votes}_L(c)/M$, was the second candidate and scored slightly lower (Section~\ref{sec:exp-graded}).

\section{Synthetic training data for the judge}\label{sec:synth}

Part~2 of the program asks whether the judge improves when synthetic data is added to its training set. The question is answered one condition at a time against the baseline C0, the judge trained on the original data. Every condition uses the same backbone and recipe and changes only the training set.

\paragraph{C1: rarity- and agreement-aware upsampling.}
C1 duplicates real records that contain rare categories or high-agreement spans, with a multiplier of up to 3, adding 8,301 records. It adds no new text and so isolates the effect of re-weighting.

\paragraph{C4: retrieval-grounded transplants (new).}
The training data has many answers per image: 12,082 answers over 2,054 images. A clean answer therefore usually has \emph{siblings}, other answers about the same image in the same language. When a sibling carries a real, annotated hallucination of the target category, C4 takes the sentence that contains it and puts it into the clean answer:
\begin{itemize}
    \item \textbf{Placement.} The sentence preferably \emph{replaces} the most similar sentence of the clean answer (token Jaccard $\ge 0.25$), so the answer still answers its question. Otherwise it is inserted at the sentence boundary closest to a position drawn from the real position distribution of that category.
    \item \textbf{Real spans.} The transplanted span is real annotated hallucination text about the same image, with its real annotator probability.
    \item \textbf{Contradiction check.} An automated check rejects placements that leave a contradicting count or a near-duplicate quoted string elsewhere in the answer, such as ``seven glasses'' next to ``five glasses''.
    \item \textbf{Scope.} C4 transplants the three rare categories and adds 719 records. \textbf{C4x} adds 1,356 transplants of the two frequent categories, with a polarity-conflict filter.
\end{itemize}

\paragraph{Realism checks.}
Before any training, each construction was compared with real spans on the training data. A classifier two-sample test~\cite{lopezpaz2017c2st} asks whether a TF-IDF logistic-regression classifier, cross-validated with folds grouped by image, can tell a method's spans, with their context, from real ones. An AUC of 0.5 means indistinguishable. Template-based heuristic injection from our shared-task system reaches 0.92--1.00, and LLM-written injections 0.87--0.91. C4 reaches 0.51--0.70 (Table~\ref{tab:c2st}).

\paragraph{Prior-matched variants (C4x-b, C1-b).}
Every added record carries a hallucination, so the conditions above also raise the share of hallucinated training records: from 71.6\% in real train to 75.7\% for C4x and 83.1\% for C1. They also change the number of optimizer steps. The prior-matched variants add clean duplicates so that the hallucinated share is exactly 71.6\%. For C4x-b the duplicates are the clean base answers the transplants were made from, so the same image and question are also seen without the hallucination. Both variants are trained for exactly C0's 3,020 optimizer steps.

\paragraph{LLM-written injections.}
A planned condition with LLM-written injections at the same locations as heuristic edits was gated on an earlier transplant condition passing a promotion rule. Since none did, it was not run and no API budget was spent (Section~\ref{sec:exp-synth}).

\section{Evaluation protocol}\label{sec:protocol}

Shared-task systems are usually tuned on one validation set many times, which overstates results~\cite{dwork2015reusable}. Every design choice in this thesis follows the protocol below, written down before the experiments and enforced in code.

\paragraph{Frozen validation halves.}
The validation set is split once, before any experiment, into \emph{val\_select} (1,515 records) and \emph{val\_confirm} (1,505 records). The split is grouped by image, so no image is on both sides, and balanced per language. Every comparison between candidates is made on val\_select only. val\_confirm only re-checks a choice that has already been made.

\paragraph{Search budget and decision log.}
A decision compares at most five candidates, globally: there are never separate choices per language or category. Each decision is appended to a log with its candidates, their scores, the choice and its confirmation, and is never reopened. Rules added during the program (R1--R6) were written down before the runs they govern.

\paragraph{Coverage guard.}
A prediction file is scored only if it covers exactly the ids of the split, and feature-based methods must show that every record had features.

\paragraph{The ``beats baseline'' criterion.}
A system $f$ beats a baseline $f_0$ trained on the same data with the same backbone if all three conditions hold:
\begin{enumerate}
    \item[(a)] its mean val\_select gain over seeds $\{42, 123, 2024\}$ exceeds the seed standard deviation of $f_0$;
    \item[(b)] the gain has the same sign on val\_confirm;
    \item[(c)] the paired bootstrap 95\% confidence interval of the val\_confirm gain lies above zero.
\end{enumerate}
The bootstrap~\cite{efron1993bootstrap} averages each record's score over the three seeds, resamples records within each language (2,000 resamples) and recomputes the difference of the four-language means. A system that meets the criterion is then submitted to the official leaderboard for all three seeds, and the test claim is the seed mean and standard deviation.
